\setgebiet{7}{Geometrie}{tief}
\chapter{Geometrie}
\marginnote{\small\itshape\color{grau}\textbf{Sie brauchen:} Band 6, Kapitel 1--14. Aus Band 2 Vektoren, Matrizen, Eigenwerte (Gebiet 2); aus Band 3 Ableitungen und Integrale (Gebiet 3); für die Kapitel 11--13 Kurven und Flächen aus Band 3, Kap.\ 13.}[12pt]
\worum{Geometrie ist das Gebiet, in dem man drei Sprachen zugleich spricht: die synthetische (Zirkel, Lineal, Kongruenz), die analytische (Koordinaten, Vektoren, Matrizen) und die differentielle (Krümmung, Geodäten). Die Aufgaben wechseln bewusst zwischen ihnen --- und zeigen, dass dieselbe Frage in jeder Sprache anders aussieht und in einer davon meist leicht ist.}

\section{Was Sie hier üben}
\begin{center}
{\footnotesize\sffamily\setlength{\tabcolsep}{4pt}
\renewcommand{\arraystretch}{1.5}
\begin{tabular}{@{}lP{4.9cm}lcll@{}}
\toprule
\textbf{Nr.} & \textbf{Titel} & \textbf{Sorte} & \textbf{Grad} & \textbf{Zeit} & \textbf{Band/Kap.}\\
\midrule
6.1 & Winkelhalbierende und Lot & \typtypisch & \grad1 & 15 min & 6/1\\
6.2 & SSW ist kein Kongruenzsatz & \typstolper & \grad1 & 10 min & 6/1\\
6.3 & Die Mittelsenkrechte & \typwerk & \grad1 & 10 min & 6/1\\
6.4 & Das 13-14-15-Dreieck & \typtypisch & \grad2 & 15 min & 6/2\\
6.5 & Die Eulergerade & \typwerk & \grad2 & 15 min & 6/2\\
6.6 & Eine Kreisgleichung lesen & \typtypisch & \grad1 & 10 min & 6/3\\
6.7 & Die Potenz eines Punktes & \typwerk & \grad1 & 10 min & 6/3\\
6.8 & Eine Ellipse vermessen & \typwerk & \grad2 & 15 min & 6/3\\
6.9 & Windschief & \typtypisch & \grad2 & 15 min & 6/4\\
6.10 & Sinus oder Kosinus? & \typstolper & \grad1 & 10 min & 6/4\\
6.11 & Ebene, Abstand, Tetraeder & \typwerk & \grad2 & 15 min & 6/4\\
6.12 & Entartete Kegelschnitte & \typstolper & \grad1 & 10 min & 6/5\\
6.13 & Typ und Normalform & \typwerk & \grad2 & 15 min & 6/5\\
6.14 & Eine Quadrik im Raum & \typwerk & \grad2 & 15 min & 6/5\\
6.15 & Drei Abbildungen, drei Typen & \typtypisch & \grad2 & 15 min & 6/6\\
6.16 & Zwei Spiegelungen & \typwerk & \grad1 & 10 min & 6/6\\
6.17 & Die Symmetriegruppe des Würfels & \typwerk & \grad2 & 15 min & 6/6\\
6.18 & Eine affine Abbildung bestimmen & \typtypisch & \grad2 & 15 min & 6/7\\
6.19 & Harmonische Lage & \typwerk & \grad1 & 10 min & 6/7\\
6.20 & Affin gibt es nur ein Dreieck & \typwerk & \grad2 & 15 min & 6/7\\
6.21 & Ein Beweis im Kreis & \typstolper & \grad2 & 15 min & 6/8\\
6.22 & Hyperbolische Vielecke & \typtypisch & \grad2 & 15 min & 6/9\\
6.23 & Ein hyperbolisches Dreieck & \typwerk & \grad1 & 10 min & 6/9\\
6.24 & Der Breitenkreis ist kein Umweg --- oder doch? & \typstolper & \grad1 & 10 min & 6/10\\
6.25 & Ein Dreieck auf der Erde & \typwerk & \grad2 & 15 min & 6/10\\
6.26 & Frankfurt--New York & \typwerk & \grad2 & 20 min & 6/10\\
6.27 & Krümmung eines Graphen & \typtypisch & \grad2 & 15 min & 6/11\\
6.28 & Die Krümmung der Ellipse & \typwerk & \grad2 & 10 min & 6/11\\
6.29 & Die Schraubenlinie & \typwerk & \grad2 & 15 min & 6/11\\
6.30 & Warum keine Landkarte stimmt & \typstolper & \grad2 & 15 min & 6/12\\
6.31 & Die Kugel und ihre Oberfläche & \typwerk & \grad2 & 15 min & 6/12\\
6.32 & Die Krümmung des Torus & \typwerk & \grad3 & 20 min & 6/12\\
6.33 & Gauß--Bonnet konkret & \typtypisch & \grad2 & 15 min & 6/13\\
6.34 & Der Igelsatz & \typwerk & \grad2 & 15 min & 6/13\\
6.35 & Die fünf platonischen Körper & \typtypisch & \grad2 & 15 min & 6/14\\
6.36 & Sylvester--Gallai & \typwerk & \grad3 & 20 min & 6/14\\
6.37 & Dichteste Packungen & \typwerk & \grad2 & 15 min & 6/14\\
\bottomrule
\end{tabular}}
\end{center}

\section{Woran Sie merken, dass Sie es können}
\begin{itemize}[leftmargin=1.2em,itemsep=2pt]
\item Sie können eine Konstruktion nicht nur ausführen, sondern mit einem Kongruenzsatz begründen --- und wissen, welche Kongruenzsätze es nicht gibt.
\item Sie wechseln frei zwischen Bild und Koordinaten: Ein Dreieck ist drei Punkte, eine Ebene ein Normalenvektor, ein Kegelschnitt eine symmetrische Matrix.
\item Sie erkennen den Typ einer Isometrie an Fixpunkten und Determinante, und Sie wissen, was affin erhalten bleibt und was nicht.
\item Sie können Winkelsumme und Fläche auf Kugel und hyperbolischer Ebene ausrechnen --- und erklären, warum beides zusammenhängt.
\end{itemize}

\section{Aufgaben}

%%%%%%%%%%% Kapitel 1
\stupser{tief}{Winkelhalbierende: gleich lange Schenkelstücke abtragen, von dort gleiche Kreise. Beweis: zwei Dreiecke mit drei gleichen Seiten.}{22pt}
\begin{karte}{tief}{6.1}{\typtypisch}{\meta{1}{15 min}{Band 6, Kap.\ 1}}
\kartentitel{Winkelhalbierende und Lot}
\worum{Zwei Grundkonstruktionen mit Zirkel und Lineal --- und für jede der Beweis, dass sie tut, was sie soll. Der Beweis ist immer ein Kongruenzsatz.}
\methoden{Konstruktion mit Zirkel und Lineal, Kongruenzsätze, Beweis einer Konstruktion}
\begin{enumerate}[label=(\alph*),leftmargin=1.6em,itemsep=2pt]
\item Konstruieren Sie die Winkelhalbierende eines Winkels mit Scheitel $S$ und beweisen Sie, dass sie den Winkel halbiert.
\item Konstruieren Sie das Lot von einem Punkt $P$ auf eine Gerade $g$ und beweisen Sie, dass es senkrecht steht.
\item Warum kann man mit Zirkel und Lineal jeden Winkel halbieren, aber nicht jeden dritteln? (Nur die Idee; der Beweis ist Band 11.)
\end{enumerate}
\end{karte}

\stupser{akzent}{Zeichnen Sie Seite $a = 5$, Winkel $\beta = 30^{\circ}$ an ihrem Ende, und einen Kreis mit Radius $b = 3$ um den anderen Endpunkt: Er schneidet den freien Schenkel zweimal.}{22pt}
\begin{karte}{akzent}{6.2}{\typstolper}{\meta{1}{10 min}{Band 6, Kap.\ 1}}
\kartentitel{SSW ist kein Kongruenzsatz}
\worum{SSS, SWS, WSW und SsW --- aber nicht SSW: Zwei Seiten und ein Winkel, der nicht der größeren Seite gegenüberliegt, bestimmen das Dreieck nicht. Wer das übersieht, beweist Dinge, die falsch sind.}
\methoden{Kongruenzsätze, Gegenbeispiel, Fallunterscheidung}
\begin{enumerate}[label=(\alph*),leftmargin=1.6em,itemsep=2pt]
\item Konstruieren Sie zwei nicht kongruente Dreiecke mit $a = 5$, $b = 3$ und $\beta = 30^{\circ}$ (Winkel gegenüber $b$).
\item Rechnen Sie mit dem Sinussatz nach: Welche Werte kann $\alpha$ haben?
\item Wann ist SSW doch eindeutig? Formulieren Sie die Bedingung (SsW).
\item Finden Sie den Fehler: \enquote{In den Dreiecken $ABC$ und $A'B'C'$ gilt $a = a'$, $b = b'$ und $\beta = \beta'$. Also sind sie kongruent, also $c = c'$.}
\end{enumerate}
\end{karte}

\stupser{gruen}{Die Schnittpunkte $P$, $Q$ der beiden Kreise haben von $A$ und $B$ denselben Abstand. Alle Punkte mit dieser Eigenschaft liegen auf einer Geraden --- warum?}{22pt}
\begin{karte}{gruen}{6.3}{\typwerk}{\meta{1}{10 min}{Band 6, Kap.\ 1, Aufgabe 1.1}}
\kartentitel{Die Mittelsenkrechte}
\worum{Die Konstruktion kennt jeder; der Beweis braucht zweimal SSS und einmal den Nebenwinkelsatz.}
\methoden{Konstruktion mit Zirkel und Lineal, Kongruenzsätze, Beweis einer Konstruktion}
Konstruieren Sie die Mittelsenkrechte einer Strecke $AB$ und beweisen Sie, dass die Konstruktion richtig ist: Die Gerade $PQ$ geht durch den Mittelpunkt und steht senkrecht.
\end{karte}

%%%%%%%%%%% Kapitel 2
\stupser{tief}{Heron: $A = \sqrt{s(s - a)(s - b)(s - c)}$ mit $s = 21$. Dann $r = A/s$, $R = abc/(4A)$.}{22pt}
\begin{karte}{tief}{6.4}{\typtypisch}{\meta{2}{15 min}{Band 6, Kap.\ 2}}
\kartentitel{Das 13-14-15-Dreieck}
\worum{Das Dreieck mit den schönsten Zahlen: Fläche, Höhen, Inkreis, Umkreis, Winkel --- alles ganz oder halb. Ein Beispiel, an dem man alle Dreiecksformeln einmal durchrechnet.}
\methoden{Heronsche Formel, Inkreis und Umkreis, Kosinussatz}
Ein Dreieck hat die Seiten $a = 13$, $b = 14$, $c = 15$.
\begin{enumerate}[label=(\alph*),leftmargin=1.6em,itemsep=2pt]
\item Berechnen Sie die Fläche mit der Heronschen Formel.
\item Berechnen Sie die Höhe auf $b$ und zerlegen Sie damit das Dreieck in zwei rechtwinklige. Was fällt auf?
\item Berechnen Sie Inkreisradius $r$ und Umkreisradius $R$.
\item Berechnen Sie den Winkel $\alpha$ mit dem Kosinussatz und prüfen Sie $A = \frac12bc\sin\alpha$.
\end{enumerate}
\end{karte}

\stupser{gruen}{Umkreismittelpunkt: Schnitt zweier Mittelsenkrechten. Höhenschnittpunkt: Schnitt zweier Höhen. Dann prüfen, ob $S$ auf der Strecke $UH$ liegt und sie $1 : 2$ teilt.}{22pt}
\begin{karte}{gruen}{6.5}{\typwerk}{\meta{2}{15 min}{Band 6, Kap.\ 2, Aufgabe 2.1}}
\kartentitel{Die Eulergerade}
\worum{Schwerpunkt, Umkreismittelpunkt und Höhenschnittpunkt liegen auf einer Geraden --- Euler 1765. Mit Koordinaten ist das eine Rechnung von zehn Zeilen.}
\methoden{Koordinatengeometrie, Besondere Punkte im Dreieck, Teilverhältnis}
Ein Dreieck hat die Ecken $A = (0, 0)$, $B = (6, 0)$, $C = (2, 4)$. Bestimmen Sie Schwerpunkt, Umkreismittelpunkt und Höhenschnittpunkt und prüfen Sie die Eulergerade.
\end{karte}

%%%%%%%%%%% Kapitel 3
\stupser{tief}{Quadratische Ergänzung in $x$ und in $y$: $(x - 3)^{2} + (y + 2)^{2} = 25$.}{22pt}
\begin{karte}{tief}{6.6}{\typtypisch}{\meta{1}{10 min}{Band 6, Kap.\ 3}}
\kartentitel{Eine Kreisgleichung lesen}
\worum{Mittelpunkt und Radius aus der ausmultiplizierten Form, eine Tangente, ein Schnitt mit einer Geraden --- Kreisrechnung, wie sie in jeder Anwendung vorkommt.}
\methoden{Quadratische Ergänzung, Kreisgleichung, Tangente}
$K\colon x^{2} + y^{2} - 6x + 4y - 12 = 0$.
\begin{enumerate}[label=(\alph*),leftmargin=1.6em,itemsep=2pt]
\item Bestimmen Sie Mittelpunkt und Radius.
\item Zeigen Sie, dass $P = (6, 2)$ auf $K$ liegt, und bestimmen Sie die Tangente in $P$.
\item Bestimmen Sie die Schnittpunkte von $K$ mit der Geraden $y = x - 1$.
\item Für welche $c$ hat $x^{2} + y^{2} - 6x + 4y + c = 0$ keinen Punkt?
\end{enumerate}
\end{karte}

\stupser{gruen}{Potenz $= d^{2} - r^{2} = 169 - 25 = 144$. Sekante: $PA\cdot PB = 144$. Tangente: $t^{2} = 144$.}{22pt}
\begin{karte}{gruen}{6.7}{\typwerk}{\meta{1}{10 min}{Band 6, Kap.\ 3, Aufgabe 3.1}}
\kartentitel{Die Potenz eines Punktes}
\worum{Sekanten-Tangenten-Satz und Sehnensatz sind ein Satz: Das Produkt der Abschnitte hängt nur vom Punkt ab, nicht von der Geraden.}
\methoden{Potenz eines Punktes, Sekanten-Tangenten-Satz}
Ein Kreis hat den Radius $5$. Ein Punkt $P$ liegt $13$ Einheiten vom Mittelpunkt entfernt. Eine Sekante durch $P$ trifft den Kreis in $A$ und $B$ mit $PA = 8$. Bestimmen Sie $PB$ und die Tangentenlänge von $P$.
\end{karte}

\stupser{gruen}{$c^{2} = a^{2} - b^{2} = 16$; $e = c/a$; Leitlinien $x = \pm a/e$. Tangente in $(x_0, y_0)$: $\frac{x_0x}{a^{2}} + \frac{y_0y}{b^{2}} = 1$.}{22pt}
\begin{karte}{gruen}{6.8}{\typwerk}{\meta{2}{15 min}{Band 6, Kap.\ 3, Aufgabe 3.2}}
\kartentitel{Eine Ellipse vermessen}
\worum{Brennpunkte, Exzentrizität, Leitlinien, Tangente --- alle Kenngrößen einer Ellipse an einem Beispiel, bei dem die Zahlen aufgehen.}
\methoden{Kegelschnitte, Brennpunkteigenschaft, Tangente}
Eine Ellipse hat die Halbachsen $a = 5$ und $b = 3$. (a) Bestimmen Sie Brennpunkte, Exzentrizität und Leitlinien. (b) Prüfen Sie für $P = (4;\ 1{,}8)$ die Brennpunkteigenschaft. (c) Bestimmen Sie die Tangente in $P$.
\end{karte}

%%%%%%%%%%% Kapitel 4
\stupser{tief}{Windschief: nicht parallel und kein Schnittpunkt. Abstand: Verbindungsvektor der Aufpunkte auf das Kreuzprodukt der Richtungen projizieren.}{22pt}
\begin{karte}{tief}{6.9}{\typtypisch}{\meta{2}{15 min}{Band 6, Kap.\ 4}}
\kartentitel{Windschief}
\worum{Zwei Geraden im Raum: parallel, schneidend oder windschief --- und im letzten Fall ein Abstand, den man mit dem Kreuzprodukt in einer Zeile bekommt.}
\methoden{Lagebeziehungen, Kreuzprodukt, Abstand}
$g\colon(1, 0, 0) + s(1, 1, 0)$, $h\colon(0, 0, 1) + t(0, 1, 2)$.
\begin{enumerate}[label=(\alph*),leftmargin=1.6em,itemsep=2pt]
\item Zeigen Sie, dass $g$ und $h$ windschief sind.
\item Berechnen Sie den Abstand von $g$ und $h$.
\item Bestimmen Sie die beiden Punkte, in denen der Abstand angenommen wird (das gemeinsame Lot).
\item Geben Sie die Ebene an, die $g$ enthält und zu $h$ parallel ist.
\end{enumerate}
\end{karte}

\stupser{akzent}{Der Winkel zwischen Gerade und Ebene ist der Winkel zur \emph{Projektion} der Geraden --- also $90^{\circ}$ minus Winkel zur Normalen. Daher $\sin$, nicht $\cos$.}{22pt}
\begin{karte}{akzent}{6.10}{\typstolper}{\meta{1}{10 min}{Band 6, Kap.\ 4}}
\kartentitel{Sinus oder Kosinus?}
\worum{Drei Winkelformeln mit dem Skalarprodukt, die man verwechselt: zwischen zwei Geraden, zwischen zwei Ebenen, zwischen Gerade und Ebene. Bei der dritten steht der Sinus --- und das übersieht jeder einmal.}
\methoden{Skalarprodukt, Winkel zwischen Gerade und Ebene, Gegenbeispiel}
\begin{enumerate}[label=(\alph*),leftmargin=1.6em,itemsep=2pt]
\item Jemand berechnet den Winkel zwischen $g\colon t(1, 1, 1)$ und der Ebene $z = 0$ als $\arccos\frac{(1, 1, 1)\cdot(0, 0, 1)}{\sqrt3} = 54{,}7^{\circ}$. Warum ist das falsch, und was ist richtig?
\item Berechnen Sie den Winkel zwischen den Ebenen $x + y = 0$ und $z = 0$. Welcher Winkel ist gemeint, wenn zwei herauskommen?
\item Berechnen Sie den Winkel zwischen $g\colon t(1, 0, 0)$ und $h\colon s(1, 1, 0)$ und zwischen $g$ und $h'\colon s(-1, -1, 0)$. Warum muss man Beträge setzen?
\item Fassen Sie die drei Formeln in einer Tabelle zusammen.
\end{enumerate}
\end{karte}

\stupser{gruen}{Normalenvektor $(B - A)\times(C - A) = (12, 7, 11)$. Abstand $= \frac{|n\cdot(D - A)|}{|n|}$, Volumen $= \frac16|n\cdot(D - A)|$.}{22pt}
\begin{karte}{gruen}{6.11}{\typwerk}{\meta{2}{15 min}{Band 6, Kap.\ 4, Aufgabe 4.1}}
\kartentitel{Ebene, Abstand, Tetraeder}
\worum{Vier Punkte im Raum, und alle Grundrechnungen daran: Ebene, Abstand, Volumen, Winkel --- mit Kreuz- und Spatprodukt.}
\methoden{Kreuzprodukt, Spatprodukt, Abstand, Skalarprodukt}
$A = (1, 2, 3)$, $B = (4, 0, 1)$, $C = (2, 5, 0)$, $D = (0, 1, 4)$. (a) Ebene durch $A, B, C$ in Koordinatenform. (b) Abstand von $D$ zu dieser Ebene. (c) Volumen des Tetraeders $ABCD$. (d) Winkel zwischen $AB$ und $AC$.
\end{karte}

%%%%%%%%%%% Kapitel 5
\stupser{akzent}{$x^{2} - y^{2} = 0$ ist $(x - y)(x + y) = 0$: zwei Geraden. $x^{2} + y^{2} = 0$: ein Punkt. $x^{2} + y^{2} = -1$: nichts.}{22pt}
\begin{karte}{akzent}{6.12}{\typstolper}{\meta{1}{10 min}{Band 6, Kap.\ 5}}
\kartentitel{Entartete Kegelschnitte}
\worum{\enquote{Jede Gleichung zweiten Grades ist Ellipse, Parabel oder Hyperbel} --- fast. Die Ausnahmen sind Geradenpaare, Punkte und die leere Menge, und man muss sie erkennen, bevor man Halbachsen ausrechnet.}
\methoden{Kegelschnitte, Fallunterscheidung, Gegenbeispiel, Faktorisieren}
Bestimmen Sie die Menge aller $(x, y)$ mit
\begin{align*} &\text{(a) } x^{2} - y^{2} = 0, &&\text{(b) } x^{2} + y^{2} + 1 = 0,\\ &\text{(c) } x^{2} - 2xy + y^{2} = 0, &&\text{(d) } x^{2} - 2xy + y^{2} - 1 = 0,\\ &\text{(e) } xy = 0 . \end{align*}
Welche sind Kegelschnitte im eigentlichen Sinn? Woran erkennt man die Entartung an der Matrix und an der Determinante?
\end{karte}

\stupser{gruen}{Matrix $\begin{pmatrix}5 & -2\\ -2 & 8\end{pmatrix}$, Eigenwerte $4$ und $9$: $4x'^{2} + 9y'^{2} = 36$.}{22pt}
\begin{karte}{gruen}{6.13}{\typwerk}{\meta{2}{15 min}{Band 6, Kap.\ 5, Aufgabe 5.1}}
\kartentitel{Typ und Normalform}
\worum{Hauptachsentransformation an einer Gleichung mit gemischtem Term: Eigenwerte bestimmen den Typ, Eigenvektoren die Achsen.}
\methoden{Hauptachsentransformation, Eigenwerte, Kegelschnitte}
Bestimmen Sie Typ und Normalform von $5x^{2} - 4xy + 8y^{2} - 36 = 0$, die Halbachsen und die Richtungen der Hauptachsen.
\end{karte}

\stupser{gruen}{$A = I + J$ mit $J$ der Einsmatrix; Eigenwerte von $J$ sind $3, 0, 0$, also die von $A$: $4, 1, 1$.}{22pt}
\begin{karte}{gruen}{6.14}{\typwerk}{\meta{2}{15 min}{Band 6, Kap.\ 5, Aufgabe 5.2}}
\kartentitel{Eine Quadrik im Raum}
\worum{Drei Variablen, alle gemischten Terme gleich --- eine Matrix mit Struktur, deren Eigenwerte man ohne charakteristisches Polynom sieht.}
\methoden{Hauptachsentransformation, Eigenwerte, Quadriken}
Klassifizieren Sie $2x^{2} + 2y^{2} + 2z^{2} + 2xy + 2yz + 2zx = 6$ und geben Sie die Halbachsen und die Symmetrieachse an.
\end{karte}

%%%%%%%%%%% Kapitel 6
\stupser{tief}{Fixpunkte suchen: $f(x) = x$. Determinante des linearen Teils: $+1$ Drehung oder Verschiebung, $-1$ Spiegelung oder Gleitspiegelung.}{22pt}
\begin{karte}{tief}{6.15}{\typtypisch}{\meta{2}{15 min}{Band 6, Kap.\ 6}}
\kartentitel{Drei Abbildungen, drei Typen}
\worum{Jede Isometrie der Ebene ist Verschiebung, Drehung, Spiegelung oder Gleitspiegelung. Die Aufgabe übt, den Typ aus der Formel abzulesen --- über Fixpunkte und Determinante.}
\methoden{Isometrien klassifizieren, Fixpunkte, Determinante}
Bestimmen Sie Typ und Kenngrößen (Zentrum, Winkel, Achse, Verschiebung) von
\begin{align*} f_1(x, y) &= (-y + 2,\ x + 1),\\ f_2(x, y) &= (y + 1,\ x + 1),\\ f_3(x, y) &= (-x + 4,\ y). \end{align*}
\end{karte}

\stupser{gruen}{Erst an $y = 0$, dann an $y = x$: Matrixprodukt $S_2S_1$. Zwei Spiegelungen an Geraden mit Winkel $\varphi$ ergeben eine Drehung um $2\varphi$.}{22pt}
\begin{karte}{gruen}{6.16}{\typwerk}{\meta{1}{10 min}{Band 6, Kap.\ 6, Aufgabe 6.1}}
\kartentitel{Zwei Spiegelungen}
\worum{Der Satz, auf dem die ganze Isometrietheorie steht: Zwei Spiegelungen sind eine Drehung (oder Verschiebung). Hier mit Matrizen, und mit der Frage nach der Reihenfolge.}
\methoden{Verkettung, Spiegelungsmatrix, Drehmatrix}
Bestimmen Sie die Verkettung der Spiegelung an $y = 0$ mit der Spiegelung an $y = x$ und identifizieren Sie den Typ. Was ändert sich bei umgekehrter Reihenfolge?
\end{karte}

\stupser{gruen}{Bahnformel: $|G| = |\text{Bahn}|\cdot|\text{Stabilisator}|$. Die Drehgruppe operiert auf $6$ Flächen, jede hat $4$ Drehungen als Stabilisator.}{22pt}
\begin{karte}{gruen}{6.17}{\typwerk}{\meta{2}{15 min}{Band 6, Kap.\ 6, Aufgabe 6.2}}
\kartentitel{Die Symmetriegruppe des Würfels}
\worum{$48$ Symmetrien, $24$ davon Drehungen --- und zwei Wege, das zu zählen: die Bahnformel und die Permutation der Raumdiagonalen.}
\methoden{Bahnformel, Symmetriegruppe, Zählen}
Bestimmen Sie die Ordnung der vollen Symmetriegruppe des Würfels auf zwei Weisen.
\end{karte}

%%%%%%%%%%% Kapitel 7
\stupser{tief}{$f(x) = Mx + t$: $t = f(0, 0)$; die Spalten von $M$ sind $f(e_1) - t$ und $f(e_2) - t$.}{22pt}
\begin{karte}{tief}{6.18}{\typtypisch}{\meta{2}{15 min}{Band 6, Kap.\ 7}}
\kartentitel{Eine affine Abbildung bestimmen}
\worum{Drei Punkte und ihre Bilder legen eine affine Abbildung fest. Was sie erhält (Parallelität, Teilverhältnisse, Flächenverhältnisse) und was nicht (Längen, Winkel), sieht man an der Matrix.}
\methoden{Affine Abbildung, Determinante, Invarianten}
$f$ sei affin mit $f(0, 0) = (1, 1)$, $f(1, 0) = (3, 2)$, $f(0, 1) = (2, 4)$.
\begin{enumerate}[label=(\alph*),leftmargin=1.6em,itemsep=2pt]
\item Bestimmen Sie $f(x) = Mx + t$.
\item Berechnen Sie das Bild des Einheitsquadrats und seine Fläche.
\item Zeigen Sie an diesem Beispiel: Mittelpunkte gehen in Mittelpunkte, parallele Geraden in parallele; rechte Winkel nicht in rechte Winkel.
\item Welche Abbildung macht das Bild wieder zum Quadrat? Ist sie affin?
\end{enumerate}
\end{karte}

\stupser{gruen}{$(A, B;\ C, D) = \frac{AC}{CB}\Big/\frac{AD}{DB} = -1$ mit orientierten Längen; $D$ liegt außerhalb der Strecke $AB$.}{22pt}
\begin{karte}{gruen}{6.19}{\typwerk}{\meta{1}{10 min}{Band 6, Kap.\ 7, Aufgabe 7.1}}
\kartentitel{Harmonische Lage}
\worum{Vier Punkte liegen harmonisch, wenn ihr Doppelverhältnis $-1$ ist: Einer teilt innen im selben Verhältnis wie der andere außen.}
\methoden{Doppelverhältnis, Teilverhältnis}
Auf einer Geraden liegen $A = 0$, $B = 6$, $C = 2$, $D = x$. Für welches $x$ liegen die vier Punkte harmonisch? Prüfen Sie, dass $C$ innen und $D$ außen im selben Verhältnis teilen.
\end{karte}

\stupser{gruen}{$u = B - A$, $v = C - A$ und entsprechend $u', v'$. Die Matrix $M$ mit $Mu = u'$, $Mv = v'$ ist $M = U'U^{-1}$.}{22pt}
\begin{karte}{gruen}{6.20}{\typwerk}{\meta{2}{15 min}{Band 6, Kap.\ 7, Aufgabe 7.2}}
\kartentitel{Affin gibt es nur ein Dreieck}
\worum{In der affinen Geometrie sind alle Dreiecke gleich --- deshalb beweist man affine Sätze am gleichseitigen oder am rechtwinkligen und hat sie für alle.}
\methoden{Affine Abbildung, Basiswechsel, Existenzbeweis}
Zeigen Sie, dass sich je zwei nichtentartete Dreiecke durch eine affine Abbildung ineinander überführen lassen. Führen Sie es durch für $A = (0, 0)$, $B = (1, 0)$, $C = (0, 1)$ und $A' = (1, 1)$, $B' = (4, 2)$, $C' = (2, 5)$. Was folgt für den Schwerpunkt?
\end{karte}

%%%%%%%%%%% Kapitel 8
\stupser{akzent}{Um zu zeigen, dass die Winkelsumme $180^{\circ}$ ist, zieht man eine Parallele durch die Spitze und benutzt Wechselwinkel. Dass diese Parallele \emph{eindeutig} ist, ist das Parallelenaxiom.}{22pt}
\begin{karte}{akzent}{6.21}{\typstolper}{\meta{2}{15 min}{Band 6, Kap.\ 8}}
\kartentitel{Ein Beweis im Kreis}
\worum{Zweitausend Jahre lang versuchte man, das Parallelenaxiom aus den anderen zu beweisen --- und benutzte dabei jedes Mal eine Aussage, die ihm äquivalent ist. Die Aufgabe übt, solche Zirkel zu erkennen.}
\methoden{Äquivalente Aussagen, Zirkelschluss erkennen, Axiomatik}
\begin{enumerate}[label=(\alph*),leftmargin=1.6em,itemsep=2pt]
\item Jemand \enquote{beweist} das Parallelenaxiom so: \enquote{Die Winkelsumme im Dreieck ist $180^{\circ}$. Gäbe es durch $P$ zwei Parallelen zu $g$, so \dots} Wo ist der Zirkel?
\item Zeigen Sie: Aus dem Parallelenaxiom folgt die Winkelsumme $180^{\circ}$. (Das ist die harmlose Richtung.)
\item Welche der folgenden Aussagen sind dem Parallelenaxiom äquivalent, welche gelten in jeder Geometrie: Es gibt ein Rechteck. Es gibt ähnliche, nicht kongruente Dreiecke. Die Winkelsumme ist höchstens $180^{\circ}$. Der Satz des Pythagoras. Durch drei nicht kollineare Punkte geht ein Kreis.
\item Warum ist die hyperbolische Geometrie der Beweis, dass es keinen Beweis gibt?
\end{enumerate}
\end{karte}

%%%%%%%%%%% Kapitel 9
\stupser{tief}{Fläche eines hyperbolischen $n$-Ecks bei $K = -1$: $(n - 2)\pi$ minus Winkelsumme. Für ein rechtwinkliges Fünfeck: $3\pi - 5\cdot\frac\pi2$.}{22pt}
\begin{karte}{tief}{6.22}{\typtypisch}{\meta{2}{15 min}{Band 6, Kap.\ 9}}
\kartentitel{Hyperbolische Vielecke}
\worum{In der hyperbolischen Ebene bestimmt die Winkelsumme die Fläche --- und es gibt Figuren, die es euklidisch nicht gibt: rechtwinklige Fünfecke, Dreiecke mit drei Nullwinkeln.}
\methoden{Winkeldefekt, Flächenformel, Gauß–Bonnet}
Krümmung $K = -1$.
\begin{enumerate}[label=(\alph*),leftmargin=1.6em,itemsep=2pt]
\item Ein Dreieck hat die Winkel $20^{\circ}$ und $30^{\circ}$ und die Fläche $1$. Bestimmen Sie den dritten Winkel.
\item Bestimmen Sie die Fläche eines Fünfecks mit fünf rechten Winkeln. Warum gibt es so ein Fünfeck euklidisch nicht?
\item Wie groß ist die Fläche eines Dreiecks höchstens, und wie sieht das Dreieck aus, das sie erreicht?
\item Kann ein hyperbolisches Dreieck die Winkel $60^{\circ}, 60^{\circ}, 60^{\circ}$ haben? Und $70^{\circ}, 70^{\circ}, 70^{\circ}$?
\end{enumerate}
\end{karte}

\stupser{gruen}{Fläche $= \pi - (\alpha + \beta + \gamma)$ bei $K = -1$; allgemein $\frac{\pi - \text{Winkelsumme}}{|K|}$.}{22pt}
\begin{karte}{gruen}{6.23}{\typwerk}{\meta{1}{10 min}{Band 6, Kap.\ 9, Aufgabe 9.1}}
\kartentitel{Ein hyperbolisches Dreieck}
\worum{Winkeldefekt gleich Fläche --- die Formel, die Gauß und Lambert fanden, an Zahlen.}
\methoden{Winkeldefekt, Flächenformel, Skalierung}
Ein hyperbolisches Dreieck hat die Winkel $30^{\circ}$, $40^{\circ}$, $50^{\circ}$. (a) Fläche bei $K = -1$. (b) Fläche bei $K = -\frac1{R^{2}}$ mit $R = 3$. (c) Welchen Anteil der maximal möglichen Dreiecksfläche nimmt es ein?
\end{karte}

%%%%%%%%%%% Kapitel 10
\stupser{akzent}{Zwei Punkte auf $60^{\circ}$ Nord, $180^{\circ}$ Länge auseinander: Entlang des Breitenkreises $\pi R\cos60^{\circ}$; über den Pol $2\cdot30^{\circ}$ Bogen.}{22pt}
\begin{karte}{akzent}{6.24}{\typstolper}{\meta{1}{10 min}{Band 6, Kap.\ 10}}
\kartentitel{Der Breitenkreis ist kein Umweg --- oder doch?}
\worum{Auf der Karte sieht der Breitenkreis wie die direkte Verbindung aus. Auf der Kugel ist er es nicht: Nur Großkreise sind Geodäten, und Breitenkreise (außer dem Äquator) sind keine.}
\methoden{Großkreis, Geodäten, Gegenbeispiel}
\begin{enumerate}[label=(\alph*),leftmargin=1.6em,itemsep=2pt]
\item Zwei Orte liegen auf $60^{\circ}$ Nord, $180^{\circ}$ Längenunterschied. Wie lang ist der Weg längs des Breitenkreises, wie lang über den Nordpol ($R = 6371$ km)?
\item Dasselbe für $10^{\circ}$ Längenunterschied auf $60^{\circ}$ Nord. Wie groß ist der Unterschied jetzt?
\item Warum ist der Äquator eine Geodäte, der Breitenkreis $60^{\circ}$ aber nicht? (Argument mit der Ebene durch den Mittelpunkt.)
\item Was bedeutet das für den Kurs eines Flugzeugs von Frankfurt nach New York (Aufgabe 6.26)?
\end{enumerate}
\end{karte}

\stupser{gruen}{Alle drei Winkel sind rechte: Exzess $\frac\pi2$, Fläche $R^{2}\cdot\frac\pi2$. Das ist ein Achtel der Kugel.}{22pt}
\begin{karte}{gruen}{6.25}{\typwerk}{\meta{2}{15 min}{Band 6, Kap.\ 10, Aufgabe 10.1}}
\kartentitel{Ein Dreieck auf der Erde}
\worum{Winkelexzess gleich Fläche durch $R^{2}$ --- das sphärische Gegenstück zum Winkeldefekt. Am Oktanten sieht man es, am $100$-km-Dreieck merkt man nichts.}
\methoden{Winkelexzess, Flächenformel, Größenordnung schätzen}
Ein sphärisches Dreieck auf der Erde ($R = 6371$ km) hat die Ecken Nordpol, $(0^{\circ}\,\text{N}, 0^{\circ}\,\text{O})$ und $(0^{\circ}\,\text{N}, 90^{\circ}\,\text{O})$. (a) Winkel und Fläche. (b) Vergleich mit dem Achtel der Kugeloberfläche. (c) Exzess eines Dreiecks mit $100$ km Seitenlänge.
\end{karte}

\stupser{gruen}{Sphärischer Kosinussatz im Dreieck Pol--Frankfurt--New York: $\cos d = \sin\varphi_1\sin\varphi_2 + \cos\varphi_1\cos\varphi_2\cos\Delta\lambda$.}{22pt}
\begin{karte}{gruen}{6.26}{\typwerk}{\meta{2}{20 min}{Band 6, Kap.\ 10, Aufgabe 10.2}}
\kartentitel{Frankfurt--New York}
\worum{Die Großkreisentfernung und der Startkurs --- sphärische Trigonometrie, wie die Navigation sie seit Jahrhunderten benutzt.}
\methoden{Sphärischer Kosinussatz, Großkreis, Sphärische Trigonometrie}
Berechnen Sie die kürzeste Flugstrecke von Frankfurt ($50{,}0^{\circ}$ N, $8{,}6^{\circ}$ O) nach New York ($40{,}7^{\circ}$ N, $74{,}0^{\circ}$ W) und den Startkurs.
\end{karte}

%%%%%%%%%%% Kapitel 11
\stupser{tief}{$\kappa = \frac{|f''|}{(1 + f'^{2})^{3/2}}$. Im Scheitel der Parabel $y = x^{2}$: $f' = 0$, $f'' = 2$.}{22pt}
\begin{karte}{tief}{6.27}{\typtypisch}{\meta{2}{15 min}{Band 6, Kap.\ 11}}
\kartentitel{Krümmung eines Graphen}
\worum{Krümmung ist nicht die zweite Ableitung --- aber dort, wo die Tangente waagerecht ist, schon. Die Formel für Graphen, Schmiegkreise, und die Bogenlänge als Kontrolle.}
\methoden{Krümmung, Schmiegkreis, Bogenlänge}
\begin{enumerate}[label=(\alph*),leftmargin=1.6em,itemsep=2pt]
\item Leiten Sie $\kappa = \frac{|f''|}{(1 + f'^{2})^{3/2}}$ aus der Krümmungsformel für Parameterkurven her.
\item Berechnen Sie die Krümmung von $y = x^{2}$ bei $x = 0$ und $x = 1$ und die Schmiegkreisradien.
\item Warum ist $\kappa\neq|f''|$, obwohl $f''$ die \enquote{Krümmung} des Graphen misst? Was misst $f''$ stattdessen?
\item Berechnen Sie die Bogenlänge der Kettenlinie $y = \cosh x$ von $0$ bis $1$.
\end{enumerate}
\end{karte}

\stupser{gruen}{$x'y'' - y'x'' = 6$ für alle $t$; der Nenner $(9\sin^{2}t + 4\cos^{2}t)^{3/2}$ ist in den Scheiteln $8$ bzw.\ $27$.}{22pt}
\begin{karte}{gruen}{6.28}{\typwerk}{\meta{2}{10 min}{Band 6, Kap.\ 11, Aufgabe 11.1}}
\kartentitel{Die Krümmung der Ellipse}
\worum{Die Ellipse ist an den Enden der großen Achse am stärksten gekrümmt --- und die Schmiegkreisradien sind $b^{2}/a$ und $a^{2}/b$.}
\methoden{Krümmung, Schmiegkreis, Parameterkurve}
Bestimmen Sie die Krümmung der Ellipse $\gamma(t) = (3\cos t, 2\sin t)$ in den Scheiteln und die Schmiegkreisradien.
\end{karte}

\stupser{gruen}{$|\gamma'| = \sqrt{a^{2} + b^{2}}$ konstant; $\kappa = \frac{|\gamma'\times\gamma''|}{|\gamma'|^{3}}$, $\tau = \frac{(\gamma'\times\gamma'')\cdot\gamma'''}{|\gamma'\times\gamma''|^{2}}$.}{22pt}
\begin{karte}{gruen}{6.29}{\typwerk}{\meta{2}{15 min}{Band 6, Kap.\ 11, Aufgabe 11.2}}
\kartentitel{Die Schraubenlinie}
\worum{Die einzige Kurve mit konstanter Krümmung \emph{und} konstanter Torsion --- und die Grenzfälle Kreis ($b = 0$) und Gerade ($a = 0$).}
\methoden{Krümmung, Torsion, Frenet-Formeln}
Bestimmen Sie Krümmung und Torsion der Schraubenlinie $\gamma(t) = (a\cos t, a\sin t, bt)$, $a > 0$. Diskutieren Sie die Grenzfälle.
\end{karte}

%%%%%%%%%%% Kapitel 12
\stupser{akzent}{$\int_0^{2\pi}\!\!\int_0^{\pi}R^{2}\,d\theta\,d\varphi = 2\pi^{2}R^{2}\neq4\pi R^{2}$: Das Flächenelement ist $R^{2}\sin\theta\,d\theta\,d\varphi$. Und: $K = \frac1{R^{2}}\neq0$ --- keine Abwicklung.}{22pt}
\begin{karte}{akzent}{6.30}{\typstolper}{\meta{2}{15 min}{Band 6, Kap.\ 12}}
\kartentitel{Warum keine Landkarte stimmt}
\worum{Zwei Irrtümer: Das Flächenelement sei $d\theta\,d\varphi$, und man könne die Kugel wie einen Zylinder in die Ebene abrollen. Beide scheitern an derselben Zahl --- der Gauß-Krümmung.}
\methoden{Flächenelement, Theorema egregium, Gegenbeispiel}
\begin{enumerate}[label=(\alph*),leftmargin=1.6em,itemsep=2pt]
\item Jemand berechnet die Kugeloberfläche als $\int_0^{2\pi}\int_0^{\pi}R^{2}\,d\theta\,d\varphi = 2\pi^{2}R^{2}$. Wo ist der Fehler?
\item Der Zylinder lässt sich ohne Verzerrung in die Ebene abwickeln, die Kugel nicht. Begründen Sie beides mit der Gauß-Krümmung.
\item Die Mercator-Karte ist winkeltreu. Warum kann sie nicht auch flächentreu sein? Was passiert mit Grönland?
\item Gibt es eine längentreue Karte eines Kugelstücks, und sei es noch so klein?
\end{enumerate}
\end{karte}

\stupser{gruen}{$X_\theta\cdot X_\theta = R^{2}$, $X_\theta\cdot X_\varphi = 0$, $X_\varphi\cdot X_\varphi = R^{2}\sin^{2}\theta$. Flächenelement $\sqrt{EG - F^{2}} = R^{2}\sin\theta$.}{22pt}
\begin{karte}{gruen}{6.31}{\typwerk}{\meta{2}{15 min}{Band 6, Kap.\ 12, Aufgabe 12.1}}
\kartentitel{Die Kugel und ihre Oberfläche}
\worum{Die erste Fundamentalform ist das Lineal auf der Fläche: Sie misst Längen, Winkel und Flächen in den Parametern.}
\methoden{Erste Fundamentalform, Flächenelement, Kugelkoordinaten}
Bestimmen Sie die erste Fundamentalform der Kugel vom Radius $R$ in Kugelkoordinaten und berechnen Sie damit die Oberfläche.
\end{karte}

\stupser{gruen}{$K = \frac{\cos u}{r(R + r\cos u)}$: Vorzeichen von $\cos u$ --- positiv außen, negativ innen, null oben und unten.}{22pt}
\begin{karte}{gruen}{6.32}{\typwerk}{\meta{3}{20 min}{Band 6, Kap.\ 12, Aufgabe 12.2}}
\kartentitel{Die Krümmung des Torus}
\worum{Eine Fläche mit allen drei Krümmungstypen: elliptisch außen, hyperbolisch innen, parabolisch dazwischen. Und das Integral über $K$ ist null --- Gauß--Bonnet für $\chi = 0$.}
\methoden{Gauß-Krümmung, Zweite Fundamentalform, Fallunterscheidung}
Berechnen Sie die Gauß-Krümmung des Torus $X(u, v) = \bigl((R + r\cos u)\cos v, (R + r\cos u)\sin v, r\sin u\bigr)$, $R > r > 0$, und bestimmen Sie, wo er elliptisch, parabolisch und hyperbolisch ist. Prüfen Sie $\int_TK\,dA = 0$.
\end{karte}

%%%%%%%%%%% Kapitel 13
\stupser{tief}{$\int_MK\,dA = 2\pi\chi(M)$. Kugel $\chi = 2$, Torus $\chi = 0$. Für Vielecke: $\int K + \text{Außenwinkel} = 2\pi$.}{22pt}
\begin{karte}{tief}{6.33}{\typtypisch}{\meta{2}{15 min}{Band 6, Kap.\ 13}}
\kartentitel{Gauß--Bonnet konkret}
\worum{Der Satz, der Krümmung (lokal) und Topologie (global) verbindet --- an Beispielen, bei denen man beide Seiten ausrechnen kann.}
\methoden{Gauß–Bonnet, Euler-Charakteristik, Winkelexzess}
\begin{enumerate}[label=(\alph*),leftmargin=1.6em,itemsep=2pt]
\item Prüfen Sie Gauß--Bonnet für die Kugel vom Radius $R$: $\int K\,dA = 2\pi\chi$.
\item Ein geodätisches Dreieck auf der Kugel mit drei rechten Winkeln: Berechnen Sie $\int K\,dA$ über das Dreieck und vergleichen Sie mit der Winkelsumme.
\item Ein Fußball hat $12$ Fünfecke und $20$ Sechsecke. Zählen Sie Ecken und Kanten und prüfen Sie $V - E + F = 2$.
\item Warum muss jede geschlossene Fläche im Raum, die kein Torus-Loch hat, irgendwo positive Krümmung haben --- egal wie verbeult sie ist?
\end{enumerate}
\end{karte}

\stupser{gruen}{Poincaré--Hopf: Die Summe der Indizes der Nullstellen eines Vektorfelds ist $\chi(M)$. Auf der Kugel $\chi = 2\neq0$, also gibt es Nullstellen.}{22pt}
\begin{karte}{gruen}{6.34}{\typwerk}{\meta{2}{15 min}{Band 6, Kap.\ 13, Aufgabe 13.1}}
\kartentitel{Der Igelsatz}
\worum{Man kann einen Igel nicht kämmen --- und auf der Erde gibt es immer einen Punkt ohne Wind. Beides ist die Euler-Charakteristik der Kugel.}
\methoden{Index eines Vektorfelds, Euler-Charakteristik, Poincaré–Hopf}
Der Igelsatz: Auf einer Kugel gibt es kein stetiges, nirgends verschwindendes Tangentialvektorfeld. Erklären Sie den Zusammenhang mit Gauß--Bonnet. Warum lässt sich der Torus kämmen?
\end{karte}

%%%%%%%%%%% Kapitel 14
\stupser{tief}{Jede Fläche hat $p$ Kanten, jede Kante gehört zu $2$ Flächen: $pF = 2E$. An jeder Ecke $q$ Kanten: $qV = 2E$. In $V - E + F = 2$ einsetzen.}{22pt}
\begin{karte}{tief}{6.35}{\typtypisch}{\meta{2}{15 min}{Band 6, Kap.\ 14}}
\kartentitel{Die fünf platonischen Körper}
\worum{Dass es genau fünf gibt, ist ein Satz --- und sein Beweis ist drei Zeilen Eulersche Polyederformel plus eine Ungleichung.}
\methoden{Eulersche Polyederformel, Doppeltes Abzählen, Ungleichung lösen}
\begin{enumerate}[label=(\alph*),leftmargin=1.6em,itemsep=2pt]
\item Ein regelmäßiges Polyeder habe $p$-eckige Flächen, $q$ an jeder Ecke. Drücken Sie $V$ und $F$ durch $E$ aus und setzen Sie in $V - E + F = 2$ ein.
\item Leiten Sie $\frac1p + \frac1q > \frac12$ her und bestimmen Sie alle Lösungen mit $p, q\ge3$.
\item Berechnen Sie für jede Lösung $V, E, F$ und benennen Sie den Körper.
\item Welche Lösungen gibt es für $\frac1p + \frac1q = \frac12$, und was beschreiben sie?
\end{enumerate}
\end{karte}

\stupser{gruen}{Unter allen Paaren (Punkt, Verbindungsgerade) mit positivem Abstand wähle das mit kleinstem Abstand. Läge auf der Geraden ein dritter Punkt, fände man ein noch kleineres Paar.}{22pt}
\begin{karte}{gruen}{6.36}{\typwerk}{\meta{3}{20 min}{Band 6, Kap.\ 14, Aufgabe 14.1}}
\kartentitel{Sylvester--Gallai}
\worum{Endlich viele Punkte, nicht alle auf einer Geraden: Dann gibt es eine Gerade durch genau zwei. Kellys Beweis (1948) ist das Musterbeispiel für das Extremalprinzip.}
\methoden{Extremalprinzip, Widerspruchsbeweis, Diskrete Geometrie}
Beweisen Sie den Satz von Sylvester--Gallai. Warum scheitert der Beweis bei unendlich vielen Punkten --- und gilt der Satz dort überhaupt?
\end{karte}

\stupser{gruen}{Ebene: Gleichseitiges Dreieck mit Seite $2r$ enthält drei Sechstelkreise, also einen halben Kreis. Dichte $= \frac{\pi r^{2}/2}{\sqrt3r^{2}}$.}{22pt}
\begin{karte}{gruen}{6.37}{\typwerk}{\meta{2}{15 min}{Band 6, Kap.\ 14, Aufgabe 14.2}}
\kartentitel{Dichteste Packungen}
\worum{$\frac\pi{2\sqrt3}\approx0{,}907$ in der Ebene (Thue 1910), $\frac\pi{3\sqrt2}\approx0{,}740$ im Raum (Kepler 1611, Hales 1998). Die Zahlen sind elementar, die Beweise nicht.}
\methoden{Elementarzelle, Dichte berechnen, Volumen}
Berechnen Sie die Dichte der dichtesten Kugelpackung in Dimension $2$ und $3$.
\end{karte}
